\documentclass[UTF8,fontset=fandol,10pt]{ctexart}
\usepackage[a4paper,margin=19mm]{geometry}
\usepackage{amsmath,amssymb,array,longtable,booktabs,url,listings,xcolor,fancyhdr}
\usepackage[colorlinks=true,linkcolor=blue,urlcolor=blue]{hyperref}
\pagestyle{fancy}\fancyhf{}\fancyhead[L]{YunTAPR-Net / candidate loss implementation}
\fancyhead[R]{2026-10-10 / v1}\fancyfoot[C]{\thepage}\setlength{\headheight}{14pt}
\setlength{\parindent}{0pt}\setlength{\parskip}{5pt}\setlength{\emergencystretch}{3em}
\lstset{language=Python,basicstyle=\ttfamily\scriptsize,columns=fullflexible,breaklines=true,
breakatwhitespace=false,numbers=left,numberstyle=\tiny,numbersep=5pt,xleftmargin=8pt,
frame=single,keepspaces=true,showstringspaces=false,tabsize=4}
\title{YunTAPR-Net\\Phase-B v2 消融候选实现\\合成验证与研究者代码学习材料}
\author{工程候选 / 独立科学审查仍待研究者}\date{2026年10月10日 / v1}
\begin{document}\maketitle
\textbf{实际实现与合成验证已完成；正式科学协议和训练仍未批准。}

CPU 177 / CUDA 3 合成测试通过，静态检查11项；没有模型forward、optimizer step、原始观测或checkpoint读取。
本轮开发委托替代上轮候选阶段的练习先行顺序，但不接受Phase-A科学结果或追认历史B1恢复。
\setcounter{tocdepth}{1}\begingroup\small\setlength{\parskip}{0pt}\tableofcontents\endgroup
\clearpage\section{候选实现架构与责任边界}

科学定义来自冻结 Focal、Pinball、total\_\allowbreak{}loss、v2 quantile 输出守卫及上轮 E0/\allowbreak{}E1/\allowbreak{}E2 提案；工程保护在隔离目录内补齐。当前委托允许候选代码和合成 autograd，未批准冻结合同修改、数据 preflight、真实模型推理或 optimizer 更新。

\begingroup\small\setlength{\tabcolsep}{3pt}\renewcommand{\arraystretch}{1.12}
\begin{longtable}{>{\raggedright\arraybackslash}p{\dimexpr(\linewidth-6\tabcolsep)/3\relax}>{\raggedright\arraybackslash}p{\dimexpr(\linewidth-6\tabcolsep)/3\relax}>{\raggedright\arraybackslash}p{\dimexpr(\linewidth-6\tabcolsep)/3\relax}}
\toprule
模块 & 输入/\allowbreak{}输出 & 责任 \\
\midrule\endfirsthead
模块 & 输入/\allowbreak{}输出 & 责任 \\
\midrule\endhead
config.py & E0/\allowbreak{}E1/\allowbreak{}E2 → 不可变参数；model/\allowbreak{}seed → RunSpec & 闭合集合、18 身份、纯数学预算与 LR 前缀 \\
validation.py & logit/\allowbreak{}qlog/\allowbreak{}rate/\allowbreak{}独立 bool masks → Supervision & 形状、设备、dtype、finite、支撑/\allowbreak{}严格序、float32 标签与计数 \\
focal.py & logit、硬标签、valid、显式 alpha/\allowbreak{}gamma → S\_\allowbreak{}occ & 稳定 BCE、类别权重、focusing、有效位置 sum \\
pinball.py & FP64 qlog/\allowbreak{}log\_\allowbreak{}target/\allowbreak{}rainy\_\allowbreak{}valid → S\_\allowbreak{}qr & 固定 32 tau、mean 后 sum，不转换物理域 \\
total.py & 两头合成张量与固定 config → CandidateLossResult & 唯一权重、N\_\allowbreak{}valid 分母、零监督停止 \\
metrics.py & 候选 result /\allowbreak{} 同输入张量 → ReportingSums & unweighted CPB、全局 pool、固定 E0 FP64 common core \\
readiness.py & 合成或声明的 SHA 元数据 → 审查结果 & fresh/\allowbreak{}恢复绑定保护；BlockedRunner 永远抛错 \\
\_\allowbreak{}\_\allowbreak{}init\_\allowbreak{}\_\allowbreak{}.py & 公开接口导入 & 没有数据或训练启动行为；导入损失 API 会加载 torch \\
\bottomrule\end{longtable}\endgroup

调用顺序：candidate\_\allowbreak{}loss → supervision → occurrence\_\allowbreak{}numerator 与 conditional\_\allowbreak{}pinball\_\allowbreak{}numerator → combine\_\allowbreak{}numerators → result。Pinball 的 \_\allowbreak{}pinball\_\allowbreak{}from\_\allowbreak{}errors 是私有算术内核，专供已经定义的 signed errors；不是绕过 qlog 资格的正式入口。P1–P4 用内核核对手算，不把常数 qlog 伪装成 v2 单调输出。

模型输出仍为 occurrence logit 与 FP64 32 条件分位数，本轮没有实现或修改 head、allocation/\allowbreak{}span 参数化、tau、epsilon、SP04 和 causal slots。validate\_\allowbreak{}log\_\allowbreak{}quantiles 是唯一正式输出守卫复用；没有调用模型构造器、Contract.load、torch.load、数据模块或旧训练入口。该复用保持冻结的 finite、q1>log1p(.1)、严格递增定义。

supervision 拒绝零有效像元而非返回 batch\_\allowbreak{}skipped；这与旧正式 runner 对 skipped batch 的最终拒绝一致，但错误提前到候选 loss 入口。有效但全无雨仍保留 occurrence 与 q 的零梯度连接。输入 rate 必须是原 float32 参考，拒绝 requires\_\allowbreak{}grad；输出全域 finite，即使像元无效也不能掩盖坏模型输出。仅 invalid reference 的临时计算副本置零，源张量不改写。

FP16/\allowbreak{}BF16 logits 显式提升为 FP32 做 BCE；FP64 logits 保持 FP64用于梯度差分。qlog/\allowbreak{}log\_\allowbreak{}target/\allowbreak{}tau 固定 FP64。autocast(enabled=False) 防止外层混合精度改写关键路径，转换的反向链仍会将梯度送回原 logits。该输入张量层测试不宣称真实 BF16 模型已通过。

静态工具不 import 候选模块，仅读源码 AST、提案和发表的元数据 SHA。readiness.py 自身只用标准库，但通过公开包导入时 \_\allowbreak{}\_\allowbreak{}init\_\allowbreak{}\_\allowbreak{} 会加载 torch；这仍没有实例化模型。真正只需标准库的入口为 scripts/\allowbreak{}.../\allowbreak{}static\_\allowbreak{}review.py。
\clearpage\section{闭合消融配置与运行身份合同}

config.py定义私有\_\allowbreak{}PARAMETERS字典：E0=(.5,2,1)，E1=(.5,0,1)，E2=(.5,2,2)。get\_\allowbreak{}config(experiment\_\allowbreak{}id)只接收精确字符串E0/\allowbreak{}E1/\allowbreak{}E2；大小写错误、未知ID、E3/\allowbreak{}E4/\allowbreak{}FinalFit、None或dict全部拒绝。它不读取外部配置文件、不选择最优条件。

AblationConfig是@dataclass(frozen=True)：自动生成构造、比较等方法，字段是experiment\_\allowbreak{}id、alpha、gamma、lambda\_\allowbreak{}q。\_\allowbreak{}\_\allowbreak{}post\_\allowbreak{}init\_\allowbreak{}\_\allowbreak{}在构造后检查类型、finite和ID对应精确数值，禁止E0名下悄悄使用gamma0。frozen禁止普通属性赋值；候选调用前仍重验，不能把Python级防篡改当人类授权。

finite\_\allowbreak{}number(name,value)返回Python float但不盲目把字符串或bool转数；判断type不是isinstance用于排除bool。无默认alpha/\allowbreak{}gamma，裸loss显式传参；科学臂只使用get\_\allowbreak{}config返回值。科学状态PROPOSED\_\allowbreak{}FOR\_\allowbreak{}RESEARCHER\_\allowbreak{}APPROVAL是常量说明，不存在升级APPROVED的setter。

RunSpec(experiment\_\allowbreak{}id,model,seed)要求模型精确B0\_\allowbreak{}MATCHED\_\allowbreak{}V2/\allowbreak{}B1\_\allowbreak{}V2，seed为真正int且属于2026/\allowbreak{}2027/\allowbreak{}2028。run\_\allowbreak{}id property动态构造条件、模型与seed身份；stage property将2026归阶段1、另两seed归阶段2，只是划分，不是执行批准。proposed\_\allowbreak{}runs依seed→arm→model顺序生成tuple18项，tuple不可追加；同一seed内条件必须共享样本、初始化和排列。

workload()只做整数ceil预算：10455/\allowbreak{}physicalbatch2→5228更新每epoch，尾batch1；9epoch47052；全18→846936。10501/\allowbreak{}valbatch8→1313，尾batch5；每epoch验证，共212706。阶段1六run282312，阶段2十二run564624。没有运行这些更新。新预算测试与旧budget\_\allowbreak{}summary核对，不重跑旧科学验收。

learning\_\allowbreak{}rate\_\allowbreak{}prefix(update)只接收1..47052的真正int；W5228、U261400，不把9epoch压缩成cosine全周期。前W步线性warmup低于1e-6可以合法；minLR只作用于后续cosine段。测试逐一核对全部47052位置与原lr\_\allowbreak{}for\_\allowbreak{}update(steps\_\allowbreak{}per\_\allowbreak{}epoch=5228)的相同返回值，原函数通过AST提取纯数学节点调用，未加载其数据/\allowbreak{}授权模块。

候选配置仍依赖D1/\allowbreak{}Q1/\allowbreak{}N0/\allowbreak{}I3/\allowbreak{}B9/\allowbreak{}S0/\allowbreak{}V0的独立正式批准。数据产品、M1配对、原float32>.1标签、32tau、epsilon1e-4、FP64守卫、batch2/\allowbreak{}accum1/\allowbreak{}no\_\allowbreak{}drop、AdamW与clip、逐epoch验证及epoch9终点不能因为代码PASS而改变。未来runner须验证机读提案、代码SHA及真实授权，而不是只检查实验ID存在。
\clearpage\section{张量契约与工程保护解读}

validation.py集中拒绝错误输入，不负责改变tau、epsilon或数学损失。ZeroValidPixelsError继承ValueError：调用者可精确捕获无监督异常，但正式runner应停止，不能catch后静默继续。Supervision为frozen dataclass，字段valid/\allowbreak{}rainy/\allowbreak{}rainy\_\allowbreak{}valid/\allowbreak{}clean\_\allowbreak{}rate都是张量，n\_\allowbreak{}valid/\allowbreak{}n\_\allowbreak{}rain是Python整数；frozen防止字段重新赋值，不会把内部Tensor变成真正只读。

image\_\allowbreak{}tensor(name,value,channels)检查torch.Tensor、torch.strided、四维[B,C,H,W]及全部维度>0。channels只接受调用方固定1或32。strided是普通dense张量布局，区别于sparse；不要求contiguous，因此合法非连续切片仍可运算。该函数不检查浮点dtype，具体入口随后检查，职责不要混淆。

same\_\allowbreak{}shape\_\allowbreak{}device(name,value,reference)要求精确同shape和device，返回None代表通过，失败抛ValueError。boolean\_\allowbreak{}mask再要求bool/\allowbreak{}dense。mask不能用float0/\allowbreak{}1替代，也不能把[H,W]悄悄扩成[B,1,H,W]；未来调用者若需构造batch掩膜必须明确产生合合同形张量。

finite\_\allowbreak{}tensor(name,value)调用torch.isfinite(value).all()再转Python bool；CUDA会同步。NaN和±Inf都拒绝。occurrence\_\allowbreak{}logits要求实浮点FP16/\allowbreak{}BF16/\allowbreak{}32/\allowbreak{}64、CPU或CUDA、单通道及全域finite；quantiles要求32通道CPU/\allowbreak{}CUDA，并复用原validate\_\allowbreak{}log\_\allowbreak{}quantiles的FP64/\allowbreak{}支撑/\allowbreak{}严格序，不自己发明epsilon或clamp。

graph\_\allowbreak{}zero(value,dtype=None)的reshape(-1)把元素拉平，[:0]取零元素切片，sum(dtype=...)产生零维0。它不先把所有值加起来，所以即使总和会溢出也能生成连接原图的0；梯度是全零。这里的零不是把非法输入修好，前置守卫仍先拒绝NaN/\allowbreak{}Inf。

supervision(logit,qlog,rate,imerg\_\allowbreak{}valid,yunnan\_\allowbreak{}mask)先检查两头与空间/\allowbreak{}device，再检查原float32非梯度rate和独立bool masks，构造交集并计数。计数0停止；有效参考负/\allowbreak{}非有限停止；invalid参考临时where成0。rainy在原精度与torch.tensor(.1,float32)比较，再与valid相交。返回的clean\_\allowbreak{}rate稍后才double()/\allowbreak{}log1p。

具体边界：float32(.1)提升后数值略高于十进制.1，所以先double再比较会错误地把边界值标有雨；测试B1保存三个nextafter相邻值。invalid参考NaN允许进入临时clean副本，但原值不改；invalid模型qlog的NaN仍拒绝，因为模型输出守卫覆盖全域。零batch、32→31通道、非boolmask、meta/\allowbreak{}CPU设备不一致都不靠广播或fallback继续。

本模块的shape/\allowbreak{}dtype/\allowbreak{}finite/\allowbreak{}异常类属于工程保护；监督交集、严格float32雨标签、q支撑/\allowbreak{}单调是冻结科学语义。未来正式网格[100,100]、3430格点及样本资格必须由另外获授权的runner检查，本loss保护不单独证明地理配准正确。
\clearpage\section{Focal 候选实现逐段解读}

阅读 config.py、validation.py 后，再打开 focal.py。occurrence\_\allowbreak{}numerator(logit,rainy,valid,*,alpha,gamma) 返回零维 S\_\allowbreak{}occ 张量，不除以 batch、有效像元数或有雨数。星号 * 表示 alpha/\allowbreak{}gamma 只能用关键字传递，减少把两参数顺序写反的风险。类型注解用于阅读；真正拒绝错误输入的是函数中的运行时检查。

logit=[B,1,H,W] 是未 sigmoid 的实数分数；rainy/\allowbreak{}valid 必须为同形、同设备的 bool。rainy 来自 rate 原 float32 严格大于 float32(.1)；Focal 裸函数不自行根据概率产生标签。valid 是 IMERG 有效性与云南掩膜交集，与 rainy 是两个不同概念。

设 b=BCEWithLogits(l,z)，pt=exp(-b)，alpha\_\allowbreak{}t=alpha*z+(1-alpha)*(1-z)，则 S\_\allowbreak{}occ=sum\_\allowbreak{}valid alpha\_\allowbreak{}t*(1-pt)\textasciicircum{}gamma*b。BCE 用 reduction='none' 保留逐元素值，再执行 focusing 和求和；不能先求 mean 再拿 scalar 计算 pt。稳定 BCE 直接吃 logits，不能先 sigmoid 再传给 BCEWithLogits。\href{https://docs.pytorch.org/docs/2.7/generated/torch.nn.functional.binary_cross_entropy_with_logits.html}{PyTorch 2.7 BCE API} 给出了输入与 reduction 的定义；项目真实数学来自 src/\allowbreak{}yuntapr/\allowbreak{}losses/\allowbreak{}focal.py。

第一个代码段调用 finite\_\allowbreak{}number：Python 中 bool 是 int 的子类，所以用 type(value) 精确排除 True/\allowbreak{}False；NaN 与 Inf 也拒绝。裸算术允许 alpha 属于 [0,1]、gamma>=0；正式候选配置只允许三个预声明组合，不开放参数搜索。随后检查 logit 与两个 masks，拒绝隐式广播。

第二段选择运算 dtype。logit.float() 的等价 .to(torch.float32) 转换不是 detach：例如原 BF16 叶子仍连接 BCE 的计算图，backward 后 grad 保留原叶子的 dtype。with torch.autocast(...,enabled=False) 只在本块内固定运算精度，退出后恢复外层上下文。

logit[valid] 是 bool 索引，将有效位置收集为 [N\_\allowbreak{}valid]；rainy[valid] 同顺序收集标签。这样 invalid 像元不进入分子，避免先做乘 mask 再指望 NaN*0 消失。这里仍先检查全域 logit finite，不能用索引隐藏坏输出。无 valid 的裸函数返回 graph\_\allowbreak{}zero；总损失入口则提前拒绝零监督。

torch.full\_\allowbreak{}like(logits,alpha) 产生同 dtype/\allowbreak{}device 的类别权重；torch.where 根据每个 hard label 选择 alpha 或 1-alpha。alpha=.5 时两类权重都为 .5，正负类别不能在 E1 改为权重 1。gamma==0 使用 alpha\_\allowbreak{}t*bce，避开不必要的 focusing 运算并保留系数；gamma2 使用 exp、pow 和 sum，与原 E0 定义一致。

F1/\allowbreak{}F3 的 l=0、gamma2 单像元期望 ln2/\allowbreak{}8；F2 的 gamma0 为 ln2/\allowbreak{}2；F4 两个相同有效元素为 ln2。103 项初轮套件包含这些手算检查；最终套件还覆盖不平衡 100 干/\allowbreak{}1 湿、±1000 和 ±1e30 logits、FP32/\allowbreak{}64及低精度输入。所有数值属于合成教学，不是科研评价结果。

gamma0 的单元素梯度为 alpha\_\allowbreak{}t*(sigmoid(l)-z)。gamma2 还需要对 focusing 权重求导，不能把 pt.detach()；测试用完整解析式和 FP64 gradcheck/\allowbreak{}gradgradcheck 对照。valid=false 的位置梯度为0。finite 检查是工程保护；alpha/\allowbreak{}gamma、BCE及 focusing 是科学定义。

小练习：独立手算 l=0,z=false,gamma0 的值和梯度；解释为什么两类 l=0 损失相同而梯度符号相反；观察 sum 与 mean 的二倍关系；找出哪些行若删掉 .5 会改变 E1。学习练习不要求重新实现才能开展本轮候选开发，但正式接入前研究者仍需理解和审查。
\clearpage\section{Conditional Pinball 逐段解读}

pinball.py 的 frozen\_\allowbreak{}taus(device=None) 永远返回 FP64 的32个值：(i-.5)/\allowbreak{}32，i=1..32，最后为 .984375。torch.arange(1,33) 的右端不包含33；先减.5再除32，不是从0到1的等间距端点集合。device 参数仅决定 CPU/\allowbreak{}CUDA位置，不改变 tau 或科学定义。

conditional\_\allowbreak{}pinball\_\allowbreak{}numerator(qlog, log\_\allowbreak{}target, rainy\_\allowbreak{}valid) 是公开入口。

qlog=[B,32,H,W]，log\_\allowbreak{}target=[B,1,H,W]，rainy\_\allowbreak{}valid=[B,1,H,W] bool；全为同设备，浮点路径必须FP64。qlog单位是 log1p(mm/\allowbreak{}h)，不是 mm/\allowbreak{}h。公开入口复用冻结 qlog finite、最低支撑和严格序检查，不允许 sort、clamp、nan\_\allowbreak{}to\_\allowbreak{}num 或 expm1。

log\_\allowbreak{}target 必须有限非负；被 rainy\_\allowbreak{}valid 选中的 log\_\allowbreak{}target 必须严格高于 log1p(float32(.1)提升至FP64)。该检查防止调用者错误地把干像元标成条件有雨。完整接口先在原 float32 构造雨标签，再提升 clean\_\allowbreak{}rate.double() 做 log1p；不允许先提升再与十进制 .1 比较。rainy\_\allowbreak{}valid=false 的高参考值可能只是被独立有效性或云南掩膜排除，不要求反向等价。

error=log\_\allowbreak{}target-qlog 的广播仅沿明确的32tau轴；不是允许 mask、batch 或空间维随意广播。e>0 表示模型分位数低于参考，惩罚 tau*e；e<0 使用(tau-1)*e，两者均非负。rho=max(tau*e,(tau-1)*e)，每个有雨像元先 mean32，再把像元损失 sum，得到未加权 S\_\allowbreak{}qr。

\_\allowbreak{}pinball\_\allowbreak{}from\_\allowbreak{}errors 是私有内核。error.movedim(1,-1) 把 [B,32,H,W] 变为 [B,H,W,32]；rainy\_\allowbreak{}valid.squeeze(1) 去掉单例通道，然后 bool 索引取出 [N\_\allowbreak{}rain,32]。tau=[32] 的乘法对应最后一轴；mean(dim=-1) 只平均tau，不平均雨像元。最后的 .sum() 不含任何 N\_\allowbreak{}valid/\allowbreak{}N\_\allowbreak{}rain 分母。

P1–P4 的常数 e=+1/\allowbreak{}-1/\allowbreak{}0 和空mask只测试内核算术：前两者 mean(tau)=mean(1-tau)=.5，第三0，第四返回连接计算图的0。这些常数 signed errors 不是常数 qlog。常数 qlog 会在公开入口被严格序守卫拒绝，不通过“关闭守卫”来运行手算测试。

无雨返回 error.reshape(-1)[:0].sum()。[:0] 是空切片，sum得到0且仍连接输入图；backward产生实际零梯度，而不是None。与 qlog.sum()*0 的数学零相同，但避免很大有限 qlog 先求和溢出后形成Inf*0。该路径属于工程实现保护，没有新增科学损失。

对e>0，dS\_\allowbreak{}qr/\allowbreak{}dq\_\allowbreak{}i=-tau\_\allowbreak{}i/\allowbreak{}32；e<0则为(1-tau\_\allowbreak{}i)/\allowbreak{}32。e=0是折点，不能用中央差分要求唯一经典导数。torch.maximum给出可用的子梯度，但确认性梯度测试刻意远离折点。\href{https://docs.pytorch.org/docs/2.7/notes/autograd.html}{PyTorch autograd说明} 可辅助理解反向图；本项目分段梯度来自实际 Pinball 定义。

本轮测试包括全部32tau解析梯度、合法严格序q、非有限输出、FP32q拒绝、干目标误选、shape/\allowbreak{}device错误、超出物理 FP64 expm1 边界的有限qlog。后者只证明 log域算术没有偷偷 materialize物理值；不证明巨大条件雨强物理可信。

小练习：对tau=.984375分别手算e=1/\allowbreak{}-1的惩罚；解释为何高tau更惩罚低估；将mean32误改sum32会使哪个量变32倍；解释覆盖率可因过宽分布提高，却不能单独证明质量改善。
\clearpage\section{总损失、监督张量与指标分母}

candidate\_\allowbreak{}loss(logit,qlog,rate,imerg\_\allowbreak{}valid,yunnan\_\allowbreak{}mask,*,config) 是候选聚合入口；config必须为闭合 AblationConfig，不能用未验证dict或临时gamma/\allowbreak{}lambda代替实验ID。再次 \_\allowbreak{}\_\allowbreak{}post\_\allowbreak{}init\_\allowbreak{}\_\allowbreak{} 校验防止外部代码绕过 frozen赋值保护后传入不匹配参数。该验证不是人类审批。

第一段 supervision：输出两头先通过形状、设备及finite/\allowbreak{}qlog守卫；参考rate必须原 float32、不参与autograd；两个mask必须独立bool且精确同形。valid=imerg\_\allowbreak{}valid \& yunnan\_\allowbreak{}mask，N\_\allowbreak{}valid=int(valid.sum().item())，rainy=clean\_\allowbreak{}rate>float32(.1)，N\_\allowbreak{}rain=sum(rainy\&valid)。.item()把零维张量转为Python计数，也会在CUDA同步，属于安全检查成本而非正式性能测量。

clean\_\allowbreak{}rate=torch.where(valid,rate,zeros\_\allowbreak{}like(rate))仅创建计算副本。有效参考NaN/\allowbreak{}Inf/\allowbreak{}负值拒绝；invalid参考允许NaN并临时置0，与原总损失一致，测试检查输入未修改。模型输出全域的NaN/\allowbreak{}Inf都拒绝，不能因invalidmask而修补。N\_\allowbreak{}valid=0抛 ZeroValidPixelsError；正常batch不静默skip，也不假称发生了更新。

第二段调用Focal得到S\_\allowbreak{}occ，再将clean\_\allowbreak{}rate提升FP64并log1p，调用Pinball得到S\_\allowbreak{}qr。发生分子已包含alpha/\allowbreak{}gamma，但“unweighted”表示没有lambda\_\allowbreak{}q；条件分子始终无lambda。第三段 combine\_\allowbreak{}numerators(S\_\allowbreak{}occ,S\_\allowbreak{}qr,N\_\allowbreak{}valid,*,lambda\_\allowbreak{}q)要求有限非负FP32/\allowbreak{}64标量、同设备、正整数N\_\allowbreak{}valid及正有限lambda，计算(S\_\allowbreak{}occ+lambda*S\_\allowbreak{}qr)/\allowbreak{}N\_\allowbreak{}valid。

示例T1：S\_\allowbreak{}occ=2,S\_\allowbreak{}qr=3,N\_\allowbreak{}valid=10，lambda1得到.5，lambda2得到.8。不是把所有loss乘2，也不是两处都乘lambda导致平方权重。T2科学CPB=3/\allowbreak{}2=1.5；N\_\allowbreak{}rain只用于科学条件评价，不用于训练分母。这些都是纯合成手算与单元测试值。

CandidateLossResult保留experiment\_\allowbreak{}id、s\_\allowbreak{}occ、s\_\allowbreak{}qr、training\_\allowbreak{}objective、weighted\_\allowbreak{}quantile\_\allowbreak{}per\_\allowbreak{}valid、n\_\allowbreak{}valid、n\_\allowbreak{}rain、conditional\_\allowbreak{}skipped、status及execution\_\allowbreak{}authorization=false。科学CPB属性只返回s\_\allowbreak{}qr/\allowbreak{}n\_\allowbreak{}rain；无雨返回None，表示不可估计，不是模型完美。weighted\_\allowbreak{}quantile\_\allowbreak{}per\_\allowbreak{}valid是辅助训练分解，不可拿它冒充unweighted CPB。

metrics.py 的 ReportingSums.from\_\allowbreak{}candidate 用detach再float生成只读汇总，停止构建梯度图但不修改训练张量。conditional\_\allowbreak{}pinball按有雨数除；pooled(other)累加分子和计数，再计算比值，不能把不等batch的CPB直接平均。common\_\allowbreak{}validation\_\allowbreak{}sums在no\_\allowbreak{}grad上下文显式使用FP64 logits及E0，固定gamma2/\allowbreak{}lambda1；训练臂的gamma或lambda不会改变共同core口径。

发生梯度和分位数梯度都通过total标量相连。gamma0时logit梯度=.5*(p-z)/\allowbreak{}N\_\allowbreak{}valid；q梯度为分段Pinball梯度再乘lambda/\allowbreak{}N\_\allowbreak{}valid。E2的q梯度为E0的2倍，而同一输入的发生梯度保持不变。这个合成张量性质不保证共享骨干真实训练后的发生输出不变。

当前candidate\_\allowbreak{}loss可以被Python调用，但没有数据读取或执行授权能力；真正runner必须在外层验证科学审批、身份、资源及执行范围。学习时只使用测试文件里的小型合成输入。本轮不生成新的真实Brier、AUROC、AP、Pinball或上尾验证结果。
\clearpage\section{未加权评价接口解读}

metrics.py没有Brier/\allowbreak{}AUROC/\allowbreak{}AP实现、没有模型推理，也没有读取样本。它只把调用者已经提供的分子与计数按冻结口径整理。正式数据是否允许传入由未来执行gate控制；本轮所有调用都是合成张量。

ReportingSums(s\_\allowbreak{}occ,s\_\allowbreak{}qr,n\_\allowbreak{}valid,n\_\allowbreak{}rain)是frozen dataclass。\_\allowbreak{}\_\allowbreak{}post\_\allowbreak{}init\_\allowbreak{}\_\allowbreak{}检查两个分子为真正Python数且finite/\allowbreak{}非负，计数是真正int、N\_\allowbreak{}valid>0、0<=N\_\allowbreak{}rain<=N\_\allowbreak{}valid；bool不接受。对象保存的是unweighted sums，不能放lambda*S\_\allowbreak{}qr进去冒充原S\_\allowbreak{}qr。接口语义由from\_\allowbreak{}candidate与测试确保，任意外部人手输入float则仍需来源审查。

from\_\allowbreak{}candidate(result)是@classmethod，首参cls代表当前类，不是对象self；它用float(result.s\_\allowbreak{}occ.detach())和s\_\allowbreak{}qr生成汇总。detach不改变数值，而是切断本报告分支的梯度图。GPU到Python scalar会同步，因此不能把汇总接口时间当纯训练吞吐。

conditional\_\allowbreak{}pinball是@property，使用者写report.conditional\_\allowbreak{}pinball而不是函数括号；返回S\_\allowbreak{}qr/\allowbreak{}N\_\allowbreak{}rain，N\_\allowbreak{}rain=0返回None。pooled(other)先加分子及计数，再创建经过同样检查的新对象，避免平均两个batch比例。它不合并seed模型预测，也不把seed\(\times\)像元当独立重复。

common\_\allowbreak{}validation\_\allowbreak{}sums(logit,qlog,rate,imerg\_\allowbreak{}valid,yunnan\_\allowbreak{}mask)被@torch.no\_\allowbreak{}grad()修饰，这些运算不建训练梯度图。它将logit.double()，显式get\_\allowbreak{}config('E0')后调用同一受保护算术入口，得到gamma2/\allowbreak{}lambda1的FP64 common core分子。训练E1的gamma0或E2的lambda2绝不能从训练result直接替换共同core。

例T2：S\_\allowbreak{}qr3、N\_\allowbreak{}rain2，CPB1.5；把这个对象与自己pooled后6/\allowbreak{}4仍1.5。若两个batch雨数不等，则应使用总S\_\allowbreak{}qr/\allowbreak{}总N\_\allowbreak{}rain，不是两个CPB简单平均。所有示例均为合成教学数值。未来指标还须绑定年份、域、样本身份、missingness及原科学审查规则，当前接口不是完整科学observer。
\clearpage\section{E0 冻结数学兼容性：实际合成验证}

证据：tests/\allowbreak{}synthetic\_\allowbreak{}attempt\_\allowbreak{}008.json与XML、tests/\allowbreak{}phase\_\allowbreak{}b\_\allowbreak{}v2\_\allowbreak{}ablation\_\allowbreak{}implementation/\allowbreak{}test\_\allowbreak{}losses.py及test\_\allowbreak{}gradients.py。最终项目环境CPU177/\allowbreak{}0、CUDA3/\allowbreak{}0；数字是检查数量，不是模型评价样本或天气事件。没有真实数据或模型推理。

\begingroup\small\setlength{\tabcolsep}{3pt}\renewcommand{\arraystretch}{1.12}
\begin{longtable}{>{\raggedright\arraybackslash}p{\dimexpr(\linewidth-8\tabcolsep)/4\relax}>{\raggedright\arraybackslash}p{\dimexpr(\linewidth-8\tabcolsep)/4\relax}>{\raggedright\arraybackslash}p{\dimexpr(\linewidth-8\tabcolsep)/4\relax}>{\raggedright\arraybackslash}p{\dimexpr(\linewidth-8\tabcolsep)/4\relax}}
\toprule
核验 & 实际范围 & 预先固定判据 & 结果 \\
\midrule\endfirsthead
核验 & 实际范围 & 预先固定判据 & 结果 \\
\midrule\endhead
Focal手算F1–F4 & 单/\allowbreak{}双元素，正负标签，gamma0/\allowbreak{}2 & FP64 rel1e-12、abs1e-14 & PASS \\
原focal函数 & FP32/\allowbreak{}FP64/\allowbreak{}BF16/\allowbreak{}FP16输入，gamma0/\allowbreak{}2 & FP32路径rtol8e-7、atol1e-8；FP64 1e-12/\allowbreak{}1e-14 & PASS \\
冻结autocast语义 & CPU BF16 logits，BCE FP32，gamma0/\allowbreak{}2 & 8e-7/\allowbreak{}1e-8 & PASS \\
原total函数与两头梯度 & B=1/\allowbreak{}2/\allowbreak{}3/\allowbreak{}5，H2/\allowbreak{}W3，独立invalid与云南mask & 8e-7/\allowbreak{}1e-8，不调宽容差 & PASS \\
未加权原Pinball & 合法严格序FP64q，雨valid交集 & 1e-12/\allowbreak{}1e-14 & PASS \\
共同validation core & 原FP64 E0分子/\allowbreak{}共同分母 & rel1e-12 & PASS \\
GPU loss与两头梯度 & 三臂，同一小型FP64合成输入 & CPU/\allowbreak{}CUDA 1e-10/\allowbreak{}1e-12 & PASS \\
标签边界 & float32(.1)及两侧nextafter & 只上邻有雨 & PASS \\
\bottomrule\end{longtable}\endgroup

相容表示合法受测输入下的数学值和梯度达到这些容差，不声明所有硬件、输入或逐位bitwise一致。候选先收集valid再sum，原函数乘mask后sum，归约顺序可能造成正常舍入差异。裸FP16/\allowbreak{}BF16在没有autocast时，原函数可能保持低精度；候选显式FP32 BCE，比较的是冻结正式autocast预期算术，而不是宣称裸低精度逐位一致。

原b0\_\allowbreak{}core\_\allowbreak{}loss没有lambda参数，E0只对应gamma2/\allowbreak{}lambda1。E1唯一改变gamma，E2唯一改变lambda；同输入科学CPB在三臂完全相同，因为q与标签不变、S\_\allowbreak{}qr无权重。未来真实训练改变共享参数后，科学结果仍未知。

有意加强的工程拒绝行为：参数NaN/\allowbreak{}Inf/\allowbreak{}bool、mask/\allowbreak{}dtype/\allowbreak{}device/\allowbreak{}shape错误、非单调或失去支撑的qlog明确拒绝；零有效监督提前抛错，旧loss先返回batch\_\allowbreak{}skipped再被正式runner拒绝。无雨仍连接q的零梯度；用空切片sum避免sum*q0的中间溢出。invalid参考沿用临时clean副本；invalid输出仍严格finite。

引用原文件字节SHA与重要函数行号见source\_\allowbreak{}identity.json；原loss、quantile参数化、正式Phase-A入口和冻结合同未修改。E0工程兼容不等于Phase-A科学验收，不追认B1历史恢复，也不能证明Brier、q32欠覆盖或极端上尾已解决。

最终边界复核的CPU007为176通过/\allowbreak{}1失败：标量设备错配应在finite算术前拒绝。代码将device检查前移并保留同形dtype合同，容差未改；CPU008=177通过/\allowbreak{}0失败，CUDA004=3通过/\allowbreak{}0失败。详见tests/\allowbreak{}scalar\_\allowbreak{}device\_\allowbreak{}repair.json与原失败日志。
\clearpage\section{合成梯度与环境验证报告}

最终项目运行时：Python3.12.14、PyTorch2.11.0+cu128、pytest9.1.1，RTX5060Laptop、CUDA12.8，构建包含sm\_\allowbreak{}120。来源tests/\allowbreak{}project\_\allowbreak{}cuda\_\allowbreak{}environment.json。测试依赖安装在本任务.local/\allowbreak{}test\_\allowbreak{}dependencies，未修改项目CUDA运行时依赖。没有测量正式训练GPU耗时、模型吞吐或显存预算。

CPU suite attempt008：177 passed、0 failures、0 errors、0 skipped；CUDA attempt004：3 passed、0 failures、0 errors、0 skipped。11项静态检查通过。当前独立测试用例180，历史attempt不相加。CPU用的是同一项目CUDA Python解释器但输入留在CPU；CUDA另开限时60秒进程。

13个CPU梯度测试：gamma0/\allowbreak{}2解析Focal梯度；gamma0/\allowbreak{}2各一次gradcheck与gradgradcheck；全部32tau的正/\allowbreak{}负分段解析梯度；E0/\allowbreak{}E1/\allowbreak{}E2的总损失双头gradcheck；三臂q梯度权重及单坐标中央差分；invalid像元零梯度。计数按pytest实例为13，gradgradcheck等多个断言仍属于其所在一个实例。

gradcheck固定FP64、eps1e-6、atol1e-6、rtol1e-5；单坐标差分eps1e-6、rel2e-8、abs1e-9。测试点远离Pinball折点和q支撑/\allowbreak{}相邻序边界。不是通过增加容差或给q排序来制造通过。Focal解析导数是alpha\_\allowbreak{}t*(p-z)*[m\textasciicircum{}gamma+gamma*b*pt*m\textasciicircum{}(gamma-1)]，m=1-pt；gamma0单独为alpha\_\allowbreak{}t*(p-z)。Pinball对q的导数在欠预测处为-tau/\allowbreak{}32，在过预测处为(1-tau)/\allowbreak{}32，再乘lambda/\allowbreak{}N\_\allowbreak{}valid。

无雨检查中q.grad是实际零张量，不是None；有效发生梯度保留。无效像元两头梯度0。P3折点仅检查值、finite与可反传，不断言唯一数学梯度。backward只求合成叶子的梯度；测试没有Optimizer实例或step，禁止torch.load/\allowbreak{}save和expm1的自动fixture会在误调用时抛错。

保留的实际尝试：attempt001的103项初轮通过；CPU002因导入冻结协议工具而缺yaml，出现1个collection error；改为只AST提取原lr\_\allowbreak{}for\_\allowbreak{}update纯函数，不import其模块。CPU003=162、004=174通过。旧torch2.7.1+cu118最小CUDA张量探测虽成功，但完整CUDA001在60秒超时，未形成测试完成数。找到历史工程记录的cu128环境后，CPU005/\allowbreak{}CUDA002因缺pytest启动失败；pytest安装到任务私有target后CPU006/\allowbreak{}CUDA003完成。相关完整输出均保留，不把未启动/\allowbreak{}超时记为通过。

这些结果仅证明受测候选数学与保护逻辑。没有真实2023/\allowbreak{}2024样本、封存2025、私有checkpoint、新模型forward或正式验证推理。核心科学假设、效应界限、多重比较和资源执行授权仍未批准。

最终边界复核的CPU007为176通过/\allowbreak{}1失败：标量设备错配应在finite算术前拒绝。代码将device检查前移并保留同形dtype合同，容差未改；CPU008=177通过/\allowbreak{}0失败，CUDA004=3通过/\allowbreak{}0失败。详见tests/\allowbreak{}scalar\_\allowbreak{}device\_\allowbreak{}repair.json与原失败日志。
\clearpage\section{独立 v2 runner 的非执行型集成提案}

当前只有readiness.BlockedRunner：describe()返回18身份、预算、阻塞事项，run(*args,**kwargs)无条件抛FormalExecutionBlocked。传入approved=true或字符串AUTHORIZED\_\allowbreak{}TO\_\allowbreak{}EXECUTE也不能运行。没有训练loop、DataLoader、模型实例、Optimizer、checkpoint应用或旧FinalFit引用。

未来建议新增隔离v2 runner，而非修改Phase-A原入口；因此本轮没有需要应用到正式路径的diff补丁。建议模块责任：审定协议/\allowbreak{}执行授权读取与独立真实性验证；数据防火墙与2023/\allowbreak{}24配对资格；可复核fresh initializer；seed和固定样本排列；训练objective适配本候选；固定E0 FP64科学汇总；逐epoch验证与immutable LAST事务；异常停止与具体LAST独立恢复；日志和资源审计。这里不实现实际训练启动。

Fresh一致性接口readiness.py逐项解释：TensorIdentity(shape,dtype,sha256)要求正整数tuple、float32、64位小写十六进制SHA。state\_\allowbreak{}metadata\_\allowbreak{}digest(tensors)将按参数名排序的声明serialize为JSON再hash，是声明层digest，不能替代真实tensor bytes哈希。FreshRecord包含RunSpec、张量映射、state/\allowbreak{}RNG/\allowbreak{}order/\allowbreak{}initializer\_\allowbreak{}code SHA及historical\_\allowbreak{}checkpoint\_\allowbreak{}loaded=false；validate重验字段与digest。

audit\_\allowbreak{}fresh\_\allowbreak{}metadata(records)要求完整18项、不重复；同模型/\allowbreak{}seed跨三臂full state digest一致；同seed六run的post\_\allowbreak{}pair RNG与样本order SHA一致；initializer代码身份一致；B0/\allowbreak{}B1名字集合相同，只允许backbone.enc0.conv1.weight和backbone.enc0.skip.weight的shape不同，其他声明完全相同。函数只检查输入metadata，不打开模型、权重或原始样本，所以actual\_\allowbreak{}initialization\_\allowbreak{}proven=false，即使metadata\_\allowbreak{}consistent=true。

后续实际initializer仍要证明参数总数、同名同形拷贝、B1两个不同形状权重保留native初始化及zero skip、独立reseed、没有historical optimizer/\allowbreak{}RNG转移、完整初态真实SHA。正式shape必须冻结501输入/\allowbreak{}100输出与模型参数量；本轮tiny metadata fixture不证明真实模型满足这些条件，也不构造这18个模型。

review\_\allowbreak{}checkpoint\_\allowbreak{}binding(metadata,expected,resume\_\allowbreak{}approval\_\allowbreak{}reference=None)要求封闭字段：run/\allowbreak{}LAST/\allowbreak{}model/\allowbreak{}optimizer/\allowbreak{}RNG/\allowbreak{}scheduler/\allowbreak{}code/\allowbreak{}protocol/\allowbreak{}data/\allowbreak{}init SHA、完成epoch、retained updates、完成train+validation边界、原执行事件reference。epoch1..9，updates=epoch*5228，expected明确绑定六项run/\allowbreak{}LAST/\allowbreak{}code/\allowbreak{}protocol/\allowbreak{}data/\allowbreak{}init。epoch9 remaining=0，不能继续训练。函数不加载或应用状态；即使给reference字符串，approval\_\allowbreak{}authenticity\_\allowbreak{}verified=false、can\_\allowbreak{}apply\_\allowbreak{}checkpoint\_\allowbreak{}state=false。

这些字符串只是合成fixture或待独立核验的来源引用，不是生成审批。未来恢复须真实研究者独立决定并绑定具体LAST SHA、原运行祖先、剩余预算、额外重放成本，再由经审查的正式gate处理；每次独立批准，禁止自动retry/\allowbreak{}resume。原B1 LAST14恢复偏差继续NOT\_\allowbreak{}GRANTED。

训练和评价必须分离：candidate\_\allowbreak{}loss按各臂gamma/\allowbreak{}lambda优化；metrics.common\_\allowbreak{}validation\_\allowbreak{}sums固定E0 FP64共同core，CPB始终unweighted S\_\allowbreak{}qr/\allowbreak{}N\_\allowbreak{}rain。18run固定epoch9，不使用新成绩选择BEST或改阈值/\allowbreak{}样本。阶段1 seed2026完整六run和阶段2另两seed十二run各有独立执行批准；阶段1成绩不能修改已批准规则或自动启动阶段2。

安全接续顺序：核对本版Git/\allowbreak{}SHA与失败记录→研究者学习并审查代码、处理科学待决字段→独立批准候选数学及正式集成范围→另版工程化新runner与获授权的模型/\allowbreak{}数据preflight→绑定真实初始化、数据、资源身份→另行批准具体阶段执行。当前开发委托已覆盖候选loss合成forward/\allowbreak{}backward，但不覆盖这些后续正式步骤。
\clearpage\section{元数据审核与阻断runner逐段解读}

readiness.py不import torch、模型、数据、optimizer，也不打开文件。其参数是Python dataclass/\allowbreak{}dict/\allowbreak{}Sequence；通过包公开入口导入仍会先执行包\_\allowbreak{}\_\allowbreak{}init\_\allowbreak{}\_\allowbreak{}并加载loss依赖，不能将整个包称为零torch依赖。static\_\allowbreak{}review.py才是标准库独立入口。

require\_\allowbreak{}sha(value,name)用re.fullmatch要求完整64位小写十六进制字符串，而非只检查某个子串。它验证格式，不证明SHA对应真实文件、更不证明人类授权。TensorIdentity(shape,dtype,sha256)的shape必须tuple且每维正整数，dtype精确float32；\_\allowbreak{}\_\allowbreak{}post\_\allowbreak{}init\_\allowbreak{}\_\allowbreak{}在构造时运行。

state\_\allowbreak{}metadata\_\allowbreak{}digest(tensors)要求非空dict，key是非空参数名，value是TensorIdentity；sorted按名字稳定排序，asdict递归把dataclass转dict，json.dumps(sort\_\allowbreak{}keys=True,separators=...)固定表示，然后hashlib.sha256(body.encode()).hexdigest()生成64位digest。该digest是声明JSON的身份，不能等同真实tensor bytes hash或加载后的全model SHA。

FreshRecord字段及validate见集成提案；validate再次检查RunSpec、四个SHA、historical\_\allowbreak{}checkpoint\_\allowbreak{}loaded严格False以及声明digest。audit\_\allowbreak{}fresh\_\allowbreak{}metadata用run\_\allowbreak{}id建查找dict，再核对完整18集合，避免重复项被dict静默覆盖。不同seed的hash没有自动相等约束；同seed跨臂和跨模型按照冻结匹配规则检查，共享shape权重须一致。

集合运算different=\{name:shape不同的名字\}在实际代码用set推导式；必须恰好两个冻结输入权重名，不允许新增第三个架构差异。其他权重比较的是TensorIdentity全字段而非只shape。这里未检查真实参数量、native skip zero或RNG来源，因为没有实际模型；输出明确actual\_\allowbreak{}initialization\_\allowbreak{}proven=false。

review\_\allowbreak{}checkpoint\_\allowbreak{}binding(metadata,expected,*,resume\_\allowbreak{}approval\_\allowbreak{}reference=None)只检查封闭字段、SHA格式、完成epoch预算与原事件reference、具体六项expected绑定。expected必须完整，不接受遗漏字段。即使reference有字符串也只是“存在待核验引用”，不会设置approval\_\allowbreak{}authenticity\_\allowbreak{}verified=true或can\_\allowbreak{}apply\_\allowbreak{}checkpoint\_\allowbreak{}state=true。epoch9无剩余更新，next\_\allowbreak{}epoch为None。

FormalExecutionBlocked继承RuntimeError。BlockedRunner.describe返回设计信息；run(*args,**kwargs)->NoReturn可以接任意伪flag，但函数唯一行为是raise。NoReturn是类型提示，不是安全机制；真正阻断来自无条件raise和没有执行分支。测试传入False/\allowbreak{}True/\allowbreak{}APPROVED/\allowbreak{}AUTHORIZED\_\allowbreak{}TO\_\allowbreak{}EXECUTE四种值都确认抛错。

练习：在合成metadata副本中改一个共享tensor SHA，预测哪个比较拒绝；删一个run或重复一项，解释为何集合检查仍能发现；把completed\_\allowbreak{}epoch改9，手算remaining=0；解释hash一致、metadata一致、真实初始化证据与独立批准为什么是四件不同事。
\clearpage\section{研究者代码与科学审查清单（未批准）}

这是待审事项，不是审批记录；没有签名、批准token或已批准勾选。当前研究者已委托候选开发与合成验证，以下真实独立审查仍需后续发生。

\noindent\textbullet\quad 阅读闭合config，确认E0/\allowbreak{}E1/\allowbreak{}E2科学假设与唯一因素；代码PASS不能替代效应界限、副作用容忍、多重比较和功效判断。\par
\noindent\textbullet\quad 解释Focal为何用logit及BCEWithLogits；gamma0为何仍是.5BCE；理解focusing梯度不是只对BCE求导。\par
\noindent\textbullet\quad 解释FP32雨标签提升前比较、独立bool masks、N\_\allowbreak{}valid/\allowbreak{}N\_\allowbreak{}rain两个分母、invalid参考临时清理与输出NaN拒绝。\par
\noindent\textbullet\quad 解释32tau、FP64 log域、mean32后sum、q支撑与严格序；区分私有signed-error手算内核和公开qlog接口。\par
\noindent\textbullet\quad 核对total仅乘一次lambda、unweighted分子保留、无雨q梯度连接、零有效明确停止；接受或提出独立候选修订，不能直接改冻结定义。\par
\noindent\textbullet\quad 阅读E0对原函数及两头梯度的实际兼容证据、固定容差、gradcheck远离折点的原因；理解合成PASS不证明真实模型或物理性能。\par
\noindent\textbullet\quad 阅读全部失败attempt，尤其旧CUDA timeout、缺yaml/\allowbreak{}pytest及任务私有依赖路径；正式GPU资源与正式依赖锁定仍未测量或批准。\par
\noindent\textbullet\quad 理解metadata fresh/\allowbreak{}恢复检查不能证明真实权重来源，也不能验证人类批准真实性；BlockedRunner无启动能力。\par
\noindent\textbullet\quad 独立决定Phase-A接受范围与B1历史恢复治理处置，保留NOT\_\allowbreak{}GRANTED；本候选不自动追认。\par
\noindent\textbullet\quad 独立批准正式集成范围后，在新版本工程化runner，开展另获授权的真实数据/\allowbreak{}模型preflight；再批准阶段1/\allowbreak{}2具体执行scope、SHA、资源和恢复规则。\par

现有2024开发验证已经用于BEST与反复诊断，不是全新确认集；q32欠覆盖、上尾异常与CPB改善未稳健确认的限制继续保留。2025、FinalFit、新架构、校准器均不纳入本候选开发的正式执行范围。
\clearpage\section{研究者代码学习手册与注释源码副本}

本轮研究者已委托Codex先实现候选，再由研究者学习与独立审查。以下副本与实际候选源码按SHA绑定，不是另一份可导入模块；修改学习副本不会修改真实实现。所有代码属于候选，最终科学决定仍由研究者作出。

\subsection{建议学习顺序与小练习}

\noindent\textbullet\quad config.py：先理解函数/\allowbreak{}关键字参数、dataclass、tuple和type检查。练习写出三臂的唯一不同项，手算18身份和每run47052更新。\par
\noindent\textbullet\quad validation.py：理解[B,C,H,W]、bool索引、shape/\allowbreak{}device/\allowbreak{}dtype与交集。练习画出两个mask交集，并解释invalid reference临时clean不等于修改源数据。\par
\noindent\textbullet\quad focal.py：对照F1–F4，先手算gamma0值/\allowbreak{}梯度，再理解gamma2对权重也求导。练习指出误删alpha会如何改变E1。\par
\noindent\textbullet\quad pinball.py：先看signed errors内核，再看公开严格q入口。练习手算高tau欠/\allowbreak{}过预测惩罚，解释mean32与sum像元的顺序。\par
\noindent\textbullet\quad total.py和metrics.py：用T1/\allowbreak{}T2核对两个分母、lambda一次与无雨None。练习合并两个不同雨数batch的CPB。\par
\noindent\textbullet\quad test\_\allowbreak{}gradients.py：认识requires\_\allowbreak{}grad、叶子、grad\_\allowbreak{}fn、backward、autograd.grad及finite difference。练习从链式法则推导q梯度的lambda/\allowbreak{}N\_\allowbreak{}valid因子。\par
\noindent\textbullet\quad readiness.py：阅读metadata而非真实state，尝试仅在个人合成副本中更改一个SHA，理解BLOCKED与数学PASS的区别。\par

\subsection{共同语法与计算图}

from .config是包内相对导入，\_\allowbreak{}\_\allowbreak{}all\_\allowbreak{}\_\allowbreak{}只定义公开名称，不授予权限。->Tensor、tuple[int,...]等是类型提示，不能替代if/\allowbreak{}raise检查。@dataclass自动生成构造与比较，frozen阻止普通字段重赋值但不冻结内部dict/\allowbreak{}Tensor。@property允许以属性形式读派生值；@classmethod的cls代表类。

Python中的and/\allowbreak{}or短路有助在错误dtype之前停止；type(x)is int排除bool。with上下文离开后恢复autocast设置。torch.where返回新Tensor；bool索引选择元素，.movedim重新安排轴，.squeeze(1)只删单例通道。reshape不等于detach，类型转换也保留可微路径。

forward是由输入算出loss并构图；backward按链式法则累加叶子.grad。autograd.grad直接返回梯度，默认不会像backward一样累加到.grad。detach/\allowbreak{}no\_\allowbreak{}grad用于报告分支，不能放到训练loss中截断梯度。gradcheck把自动微分与FP64中央差分比较；在折点不适用唯一经典导数。当前仅合成叶子参与图，没有模型参数更新或Optimizer。

工程保护包括finite/\allowbreak{}shape/\allowbreak{}device/\allowbreak{}type拒绝、SHA检查和BlockedRunner；科学定义包括alpha/\allowbreak{}gamma、32tau、log域、雨阈值和分母。两者都重要，但工程PASS不能自动批准科学假设。\href{https://docs.pytorch.org/docs/2.7/notes/autograd.html}{PyTorch autograd}及\href{https://docs.pytorch.org/docs/2.7/amp.html}{混合精度说明}提供API背景；项目公式与来源定位见对应walkthrough/\allowbreak{}source\_\allowbreak{}identity。

\subsection{测试和交付工具也应理解}

conftest.py的autouse fixture在每个测试前禁止step、torch.load/\allowbreak{}save和expm1，结束后由monkeypatch恢复原方法。synthetic\_\allowbreak{}case只用linspace/\allowbreak{}arange/\allowbreak{}tensor构造小张量，clone().requires\_\allowbreak{}grad\_\allowbreak{}()建立可观察梯度的叶子。rate/\allowbreak{}masks没有梯度。test\_\allowbreak{}losses用手算和冻结函数，不与新实现自己互证；test\_\allowbreak{}gradients用解析式/\allowbreak{}差分；test\_\allowbreak{}readiness只改内存metadata；test\_\allowbreak{}cuda是可选独立进程的三臂CPU/\allowbreak{}GPU对照。

run\_\allowbreak{}checks.py调用pytest子进程而不是模型runner，60秒timeout会终止自己启动的测试进程，保存完整输出与XML，不中断别人的任务。static\_\allowbreak{}review.py只AST与SHA，无torch import；author\_\allowbreak{}delivery.py/\allowbreak{}后续export/\allowbreak{}close工具处理文档和元数据，不读观测或权重。若修改真实候选源码，必须重新测试并用新版本身份归档，不能只改副本或覆盖已发表结果。

\subsection{\_\allowbreak{}\_\allowbreak{}init\_\allowbreak{}\_\allowbreak{}.py：实际源码、函数位置与副本}

来源 src/\allowbreak{}yuntapr/\allowbreak{}experimental/\allowbreak{}phase\_\allowbreak{}b\_\allowbreak{}v2\_\allowbreak{}ablations/\allowbreak{}\_\allowbreak{}\_\allowbreak{}init\_\allowbreak{}\_\allowbreak{}.py。

SHA256：55207122a41d\allowbreak{}d3112f393375\allowbreak{}0f78a1e23182\allowbreak{}49a38b01f37f\allowbreak{}4d00badec3de\allowbreak{}418b。参数语义见对应walkthrough；下面是实际完整注释源码副本。

\begingroup\small\setlength{\tabcolsep}{3pt}\renewcommand{\arraystretch}{1.12}
\begin{longtable}{>{\raggedright\arraybackslash}p{\dimexpr(\linewidth-6\tabcolsep)/3\relax}>{\raggedright\arraybackslash}p{\dimexpr(\linewidth-6\tabcolsep)/3\relax}>{\raggedright\arraybackslash}p{\dimexpr(\linewidth-6\tabcolsep)/3\relax}}
\toprule
函数或方法（参数名） & 源码行 & 阅读入口 \\
\midrule\endfirsthead
函数或方法（参数名） & 源码行 & 阅读入口 \\
\midrule\endhead
\bottomrule\end{longtable}\endgroup

\begin{lstlisting}
"""Candidate loss APIs only: no model, data loader, optimizer or training loop."""
from .config import AblationConfig, RunSpec, get_config, proposed_runs
from .focal import occurrence_numerator
from .pinball import conditional_pinball_numerator
from .total import CandidateLossResult, candidate_loss, combine_numerators

__all__ = ["AblationConfig", "RunSpec", "get_config", "proposed_runs",
           "occurrence_numerator", "conditional_pinball_numerator",
           "CandidateLossResult", "candidate_loss", "combine_numerators"]
\end{lstlisting}


\subsection{config.py：实际源码、函数位置与副本}

来源 src/\allowbreak{}yuntapr/\allowbreak{}experimental/\allowbreak{}phase\_\allowbreak{}b\_\allowbreak{}v2\_\allowbreak{}ablations/\allowbreak{}config.py。

SHA256：5ae850c28417\allowbreak{}b555dc7825bd\allowbreak{}204f95f8275b\allowbreak{}c3ddc5b64cbf\allowbreak{}a230a2517d4d\allowbreak{}5e3a。参数语义见对应walkthrough；下面是实际完整注释源码副本。

\begingroup\small\setlength{\tabcolsep}{3pt}\renewcommand{\arraystretch}{1.12}
\begin{longtable}{>{\raggedright\arraybackslash}p{\dimexpr(\linewidth-6\tabcolsep)/3\relax}>{\raggedright\arraybackslash}p{\dimexpr(\linewidth-6\tabcolsep)/3\relax}>{\raggedright\arraybackslash}p{\dimexpr(\linewidth-6\tabcolsep)/3\relax}}
\toprule
函数或方法（参数名） & 源码行 & 阅读入口 \\
\midrule\endfirsthead
函数或方法（参数名） & 源码行 & 阅读入口 \\
\midrule\endhead
finite\_\allowbreak{}number(name, value) & 12 & 完整type/\allowbreak{}default见下面源码 \\
get\_\allowbreak{}config(experiment\_\allowbreak{}id) & 34 & 完整type/\allowbreak{}default见下面源码 \\
proposed\_\allowbreak{}runs() & 62 & 完整type/\allowbreak{}default见下面源码 \\
workload() & 66 & 完整type/\allowbreak{}default见下面源码 \\
learning\_\allowbreak{}rate\_\allowbreak{}prefix(update) & 75 & 完整type/\allowbreak{}default见下面源码 \\
\_\allowbreak{}\_\allowbreak{}post\_\allowbreak{}init\_\allowbreak{}\_\allowbreak{}(self) & 26 & 完整type/\allowbreak{}default见下面源码 \\
\_\allowbreak{}\_\allowbreak{}post\_\allowbreak{}init\_\allowbreak{}\_\allowbreak{}(self) & 46 & 完整type/\allowbreak{}default见下面源码 \\
run\_\allowbreak{}id(self) & 54 & 完整type/\allowbreak{}default见下面源码 \\
stage(self) & 58 & 完整type/\allowbreak{}default见下面源码 \\
\bottomrule\end{longtable}\endgroup

\begin{lstlisting}
"""Three closed, unapproved loss conditions and 18 candidate run identities."""
from __future__ import annotations
from dataclasses import dataclass
import math

_PARAMETERS = {"E0": (.5, 2., 1.), "E1": (.5, 0., 1.), "E2": (.5, 2., 2.)}
MODELS = ("B0_MATCHED_V2", "B1_V2")
SEEDS = (2026, 2027, 2028)
PROTOCOL_STATUS = "PROPOSED_FOR_RESEARCHER_APPROVAL"


def finite_number(name: str, value: object) -> float:
    """Reject booleans, strings, NaN and infinity rather than silently coerce."""
    if type(value) not in (int, float) or not math.isfinite(value):
        raise ValueError(f"{name}: finite Python int/float required, excluding bool")
    return float(value)


@dataclass(frozen=True)
class AblationConfig:
    experiment_id: str
    alpha: float
    gamma: float
    lambda_q: float

    def __post_init__(self) -> None:
        if type(self.experiment_id) is not str or self.experiment_id not in _PARAMETERS:
            raise ValueError("Unknown experiment ID; only E0/E1/E2 are candidates")
        values = tuple(finite_number(k, getattr(self, k)) for k in ("alpha", "gamma", "lambda_q"))
        if values != _PARAMETERS[self.experiment_id]:
            raise ValueError("Experiment parameters do not match the closed candidate condition")


def get_config(experiment_id: str) -> AblationConfig:
    if type(experiment_id) is not str or experiment_id not in _PARAMETERS:
        raise ValueError("Unknown experiment ID; no parameter search or automatic selection")
    return AblationConfig(experiment_id, *_PARAMETERS[experiment_id])


@dataclass(frozen=True)
class RunSpec:
    experiment_id: str
    model: str
    seed: int

    def __post_init__(self) -> None:
        get_config(self.experiment_id)
        if type(self.model) is not str or self.model not in MODELS:
            raise ValueError("Unknown candidate model identity")
        if type(self.seed) is not int or self.seed not in SEEDS:
            raise ValueError("Candidate seeds are exactly 2026/2027/2028")

    @property
    def run_id(self) -> str:
        return f"{self.experiment_id}__{self.model}__s{self.seed}"

    @property
    def stage(self) -> int:
        return 1 if self.seed == 2026 else 2


def proposed_runs() -> tuple[RunSpec, ...]:
    return tuple(RunSpec(arm, model, seed) for seed in SEEDS for arm in _PARAMETERS for model in MODELS)


def workload() -> dict[str, int]:
    """Pure arithmetic; these counts do not mean that any updates took place."""
    train_batches, val_batches = (10455 + 1) // 2, (10501 + 7) // 8
    return {"runs": 18, "epochs": 9, "updates_per_epoch": train_batches,
            "updates_per_run": 9 * train_batches, "updates_total": 18 * 9 * train_batches,
            "validation_batches_per_epoch": val_batches, "validation_batches_total": 18 * 9 * val_batches,
            "stage1_updates": 6 * 9 * train_batches, "stage2_updates": 12 * 9 * train_batches}


def learning_rate_prefix(update: int) -> float:
    """First nine epochs of the original 50-epoch trajectory, no optimizer."""
    if type(update) is not int or not 1 <= update <= 47052:
        raise ValueError("Candidate one-based update must be in [1,47052]")
    if update <= 5228:
        return 1e-4 * (update / 5228)  # warmup is not clamped to the cosine floor
    progress = (update - 5228) / (261400 - 5228)
    return 1e-6 + (1e-4 - 1e-6) * (1 + math.cos(math.pi * progress)) / 2
\end{lstlisting}


\subsection{validation.py：实际源码、函数位置与副本}

来源 src/\allowbreak{}yuntapr/\allowbreak{}experimental/\allowbreak{}phase\_\allowbreak{}b\_\allowbreak{}v2\_\allowbreak{}ablations/\allowbreak{}validation.py。

SHA256：05885c309b6c\allowbreak{}f930a57f30b2\allowbreak{}5438c8664be2\allowbreak{}ddc5a36eb9f8\allowbreak{}6de82654e886\allowbreak{}beb0。参数语义见对应walkthrough；下面是实际完整注释源码副本。

\begingroup\small\setlength{\tabcolsep}{3pt}\renewcommand{\arraystretch}{1.12}
\begin{longtable}{>{\raggedright\arraybackslash}p{\dimexpr(\linewidth-6\tabcolsep)/3\relax}>{\raggedright\arraybackslash}p{\dimexpr(\linewidth-6\tabcolsep)/3\relax}>{\raggedright\arraybackslash}p{\dimexpr(\linewidth-6\tabcolsep)/3\relax}}
\toprule
函数或方法（参数名） & 源码行 & 阅读入口 \\
\midrule\endfirsthead
函数或方法（参数名） & 源码行 & 阅读入口 \\
\midrule\endhead
image\_\allowbreak{}tensor(name, value, channels) & 12 & 完整type/\allowbreak{}default见下面源码 \\
same\_\allowbreak{}shape\_\allowbreak{}device(name, value, reference) & 19 & 完整type/\allowbreak{}default见下面源码 \\
boolean\_\allowbreak{}mask(name, value, reference) & 26 & 完整type/\allowbreak{}default见下面源码 \\
finite\_\allowbreak{}tensor(name, value) & 32 & 完整type/\allowbreak{}default见下面源码 \\
occurrence\_\allowbreak{}logits(logit) & 37 & 完整type/\allowbreak{}default见下面源码 \\
quantiles(qlog) & 46 & 完整type/\allowbreak{}default见下面源码 \\
graph\_\allowbreak{}zero(value, *, dtype) & 54 & 完整type/\allowbreak{}default见下面源码 \\
supervision(logit, qlog, rate, imerg\_\allowbreak{}valid, yunnan\_\allowbreak{}mask) & 69 & 完整type/\allowbreak{}default见下面源码 \\
\bottomrule\end{longtable}\endgroup

\begin{lstlisting}
"""Explicit tensor contracts; no output repair and no physical conversion."""
from __future__ import annotations
from dataclasses import dataclass
import torch
from yuntapr.models.quantile_v2.outputs import validate_log_quantiles


class ZeroValidPixelsError(ValueError):
    """Stop a loss call with zero supervision; never imply a skipped update."""


def image_tensor(name: str, value: torch.Tensor, channels: int) -> None:
    if not isinstance(value, torch.Tensor) or value.layout != torch.strided:
        raise ValueError(f"{name}: dense strided Tensor required")
    if value.ndim != 4 or value.shape[1] != channels or any(n <= 0 for n in value.shape):
        raise ValueError(f"{name}: nonempty [B,{channels},H,W] required")


def same_shape_device(name: str, value: torch.Tensor, reference: torch.Tensor) -> None:
    if not isinstance(value, torch.Tensor) or value.shape != reference.shape:
        raise ValueError(f"{name}: exact shape required; mask broadcasting is forbidden")
    if value.device != reference.device:
        raise ValueError(f"{name}: device mismatch")


def boolean_mask(name: str, value: torch.Tensor, reference: torch.Tensor) -> None:
    same_shape_device(name, value, reference)
    if value.dtype != torch.bool or value.layout != torch.strided:
        raise ValueError(f"{name}: dense bool Tensor required")


def finite_tensor(name: str, value: torch.Tensor) -> None:
    if not bool(torch.isfinite(value).all()):
        raise FloatingPointError(f"{name}: NaN/Inf forbidden, including masked output pixels")


def occurrence_logits(logit: torch.Tensor) -> None:
    image_tensor("logit", logit, 1)
    if logit.dtype not in (torch.float16, torch.bfloat16, torch.float32, torch.float64):
        raise ValueError("logit: real floating dtype required")
    if logit.device.type not in ("cpu", "cuda"):
        raise ValueError("logit: this candidate supports CPU/CUDA only")
    finite_tensor("logit", logit)


def quantiles(qlog: torch.Tensor) -> None:
    image_tensor("qlog", qlog, 32)
    if qlog.device.type not in ("cpu", "cuda"):
        raise ValueError("qlog: this candidate supports CPU/CUDA only")
    # Reuse the frozen FP64, finite, support and strict-order guards unchanged.
    validate_log_quantiles(qlog)


def graph_zero(value: torch.Tensor, *, dtype: torch.dtype | None = None) -> torch.Tensor:
    """Empty-view sum gives connected zero without overflowing sum(value)*0."""
    return value.reshape(-1)[:0].sum(dtype=dtype)


@dataclass(frozen=True)
class Supervision:
    valid: torch.Tensor
    rainy: torch.Tensor
    rainy_valid: torch.Tensor
    clean_rate: torch.Tensor
    n_valid: int
    n_rain: int


def supervision(logit: torch.Tensor, qlog: torch.Tensor, rate: torch.Tensor,
                imerg_valid: torch.Tensor, yunnan_mask: torch.Tensor) -> Supervision:
    occurrence_logits(logit)
    same_shape_device("qlog spatial slice", qlog[:, :1] if isinstance(qlog, torch.Tensor) and qlog.ndim == 4 else qlog, logit)
    quantiles(qlog)
    same_shape_device("rate", rate, logit)
    if rate.dtype != torch.float32 or rate.layout != torch.strided or rate.requires_grad:
        raise ValueError("rate: nondifferentiable dense float32 IMERG reference required")
    boolean_mask("imerg_valid", imerg_valid, logit)
    boolean_mask("yunnan_mask", yunnan_mask, logit)
    valid = imerg_valid & yunnan_mask
    n_valid = int(valid.sum().item())
    if n_valid == 0:
        raise ZeroValidPixelsError("ZERO_VALID_SUPERVISED_YUNNAN_PIXELS; no silent skip")
    if not bool(torch.isfinite(rate[valid]).all()) or bool((rate[valid] < 0).any()):
        raise ValueError("Supervised valid precipitation must be finite nonnegative")
    # Only an ephemeral arithmetic target is cleaned, exactly as in the old loss.
    clean = torch.where(valid, rate, torch.zeros_like(rate))
    rainy = clean > torch.tensor(.1, dtype=torch.float32, device=rate.device)
    rainy_valid = rainy & valid  # compare BEFORE converting rate to FP64
    return Supervision(valid, rainy, rainy_valid, clean, n_valid, int(rainy_valid.sum().item()))
\end{lstlisting}


\subsection{focal.py：实际源码、函数位置与副本}

来源 src/\allowbreak{}yuntapr/\allowbreak{}experimental/\allowbreak{}phase\_\allowbreak{}b\_\allowbreak{}v2\_\allowbreak{}ablations/\allowbreak{}focal.py。

SHA256：828189dfac57\allowbreak{}fdaf0520ffc8\allowbreak{}a37fb4b7608e\allowbreak{}ccbd6154f3c0\allowbreak{}4538049fc531\allowbreak{}3f94。参数语义见对应walkthrough；下面是实际完整注释源码副本。

\begingroup\small\setlength{\tabcolsep}{3pt}\renewcommand{\arraystretch}{1.12}
\begin{longtable}{>{\raggedright\arraybackslash}p{\dimexpr(\linewidth-6\tabcolsep)/3\relax}>{\raggedright\arraybackslash}p{\dimexpr(\linewidth-6\tabcolsep)/3\relax}>{\raggedright\arraybackslash}p{\dimexpr(\linewidth-6\tabcolsep)/3\relax}}
\toprule
函数或方法（参数名） & 源码行 & 阅读入口 \\
\midrule\endfirsthead
函数或方法（参数名） & 源码行 & 阅读入口 \\
\midrule\endhead
occurrence\_\allowbreak{}numerator(logit, rainy, valid, *, alpha, gamma) & 9 & 完整type/\allowbreak{}default见下面源码 \\
\bottomrule\end{longtable}\endgroup

\begin{lstlisting}
"""Stable candidate focal numerator, explicitly separate from normalization."""
from __future__ import annotations
import torch
from torch.nn import functional as F
from .config import finite_number
from .validation import boolean_mask, finite_tensor, graph_zero, occurrence_logits


def occurrence_numerator(logit: torch.Tensor, rainy: torch.Tensor, valid: torch.Tensor,
                         *, alpha: float, gamma: float) -> torch.Tensor:
    alpha, gamma = finite_number("alpha", alpha), finite_number("gamma", gamma)
    if not 0 <= alpha <= 1 or gamma < 0:
        raise ValueError("alpha in [0,1] and gamma>=0 required")
    occurrence_logits(logit)
    boolean_mask("rainy hard labels", rainy, logit)
    boolean_mask("valid", valid, logit)
    dtype = torch.float64 if logit.dtype == torch.float64 else torch.float32
    with torch.autocast(logit.device.type, enabled=False):
        if not bool(valid.any()):
            return graph_zero(logit, dtype=dtype)
        # FP16/BF16 inputs retain gradients through this cast; BCE computes FP32.
        logits, labels = logit[valid].to(dtype), rainy[valid]
        bce = F.binary_cross_entropy_with_logits(logits, labels.to(dtype), reduction="none")
        alpha_t = torch.where(labels, torch.full_like(logits, alpha), torch.full_like(logits, 1 - alpha))
        if gamma == 0:
            terms = alpha_t * bce  # deliberately .5*BCE for E1, not plain BCE
        else:
            pt = torch.exp(-bce)
            terms = alpha_t * (1 - pt).pow(gamma) * bce
        result = terms.sum()
    finite_tensor("occurrence numerator", result)
    return result
\end{lstlisting}


\subsection{pinball.py：实际源码、函数位置与副本}

来源 src/\allowbreak{}yuntapr/\allowbreak{}experimental/\allowbreak{}phase\_\allowbreak{}b\_\allowbreak{}v2\_\allowbreak{}ablations/\allowbreak{}pinball.py。

SHA256：cc4c4c742a2c\allowbreak{}e5c465dc8c2e\allowbreak{}ee786b85378d\allowbreak{}4d8845d7e065\allowbreak{}783f986bc264\allowbreak{}e25f。参数语义见对应walkthrough；下面是实际完整注释源码副本。

\begingroup\small\setlength{\tabcolsep}{3pt}\renewcommand{\arraystretch}{1.12}
\begin{longtable}{>{\raggedright\arraybackslash}p{\dimexpr(\linewidth-6\tabcolsep)/3\relax}>{\raggedright\arraybackslash}p{\dimexpr(\linewidth-6\tabcolsep)/3\relax}>{\raggedright\arraybackslash}p{\dimexpr(\linewidth-6\tabcolsep)/3\relax}}
\toprule
函数或方法（参数名） & 源码行 & 阅读入口 \\
\midrule\endfirsthead
函数或方法（参数名） & 源码行 & 阅读入口 \\
\midrule\endhead
frozen\_\allowbreak{}taus(device) & 7 & 完整type/\allowbreak{}default见下面源码 \\
\_\allowbreak{}pinball\_\allowbreak{}from\_\allowbreak{}errors(error, rainy\_\allowbreak{}valid) & 11 & 完整type/\allowbreak{}default见下面源码 \\
conditional\_\allowbreak{}pinball\_\allowbreak{}numerator(qlog, log\_\allowbreak{}target, rainy\_\allowbreak{}valid) & 32 & 完整type/\allowbreak{}default见下面源码 \\
\bottomrule\end{longtable}\endgroup

\begin{lstlisting}
"""FP64 conditional pinball in log1p(mm/h), fixed 32-tau mean then sum."""
from __future__ import annotations
import torch
from .validation import boolean_mask, finite_tensor, graph_zero, image_tensor, quantiles, same_shape_device


def frozen_taus(device: torch.device | str | None = None) -> torch.Tensor:
    return (torch.arange(1, 33, device=device, dtype=torch.float64) - .5) / 32


def _pinball_from_errors(error: torch.Tensor, rainy_valid: torch.Tensor) -> torch.Tensor:
    """Private arithmetic kernel: signed errors are NOT model quantiles.

    This permits the P1-P4 hand-calculated constant-error fixtures without
    pretending that constant qlog is a legal strictly ordered v2 prediction.
    """
    image_tensor("error", error, 32)
    if error.dtype != torch.float64:
        raise ValueError("Pinball signed errors require FP64")
    boolean_mask("rainy_valid", rainy_valid, error[:, :1])
    finite_tensor("pinball errors", error)
    if not bool(rainy_valid.any()):
        return graph_zero(error)
    selected = error.movedim(1, -1)[rainy_valid.squeeze(1)]  # [N_rain,32]
    tau = frozen_taus(error.device)
    per_pixel = torch.maximum(tau * selected, (tau - 1) * selected).mean(dim=-1)
    result = per_pixel.sum()
    finite_tensor("quantile numerator", result)
    return result


def conditional_pinball_numerator(qlog: torch.Tensor, log_target: torch.Tensor,
                                  rainy_valid: torch.Tensor) -> torch.Tensor:
    quantiles(qlog)
    same_shape_device("log_target", log_target, qlog[:, :1])
    image_tensor("log_target", log_target, 1)
    if log_target.dtype != torch.float64:
        raise ValueError("log_target requires FP64")
    boolean_mask("rainy_valid", rainy_valid, log_target)
    finite_tensor("log_target", log_target)
    if bool((log_target < 0).any()):
        raise ValueError("log_target must be nonnegative log1p(mm/h)")
    threshold_log = torch.log1p(torch.tensor(.1, dtype=torch.float32, device=qlog.device).double())
    if bool((log_target[rainy_valid] <= threshold_log).any()):
        raise ValueError("rainy_valid cannot select a dry or threshold-equal reference target")
    with torch.autocast(qlog.device.type, enabled=False):
        # Only the explicitly declared singleton tau axis is expanded.
        error = log_target - qlog
        return _pinball_from_errors(error, rainy_valid)
\end{lstlisting}


\subsection{total.py：实际源码、函数位置与副本}

来源 src/\allowbreak{}yuntapr/\allowbreak{}experimental/\allowbreak{}phase\_\allowbreak{}b\_\allowbreak{}v2\_\allowbreak{}ablations/\allowbreak{}total.py。

SHA256：9ce37393d2af\allowbreak{}2fbc762d198b\allowbreak{}2088cb4645d1\allowbreak{}04d35f81b507\allowbreak{}c236bc3b39cf\allowbreak{}1ec1。参数语义见对应walkthrough；下面是实际完整注释源码副本。

\begingroup\small\setlength{\tabcolsep}{3pt}\renewcommand{\arraystretch}{1.12}
\begin{longtable}{>{\raggedright\arraybackslash}p{\dimexpr(\linewidth-6\tabcolsep)/3\relax}>{\raggedright\arraybackslash}p{\dimexpr(\linewidth-6\tabcolsep)/3\relax}>{\raggedright\arraybackslash}p{\dimexpr(\linewidth-6\tabcolsep)/3\relax}}
\toprule
函数或方法（参数名） & 源码行 & 阅读入口 \\
\midrule\endfirsthead
函数或方法（参数名） & 源码行 & 阅读入口 \\
\midrule\endhead
combine\_\allowbreak{}numerators(s\_\allowbreak{}occ, s\_\allowbreak{}qr, n\_\allowbreak{}valid, *, lambda\_\allowbreak{}q) & 11 & 完整type/\allowbreak{}default见下面源码 \\
candidate\_\allowbreak{}loss(logit, qlog, rate, imerg\_\allowbreak{}valid, yunnan\_\allowbreak{}mask, *, config) & 51 & 完整type/\allowbreak{}default见下面源码 \\
scientific\_\allowbreak{}conditional\_\allowbreak{}pinball(self) & 46 & 完整type/\allowbreak{}default见下面源码 \\
\bottomrule\end{longtable}\endgroup

\begin{lstlisting}
"""Candidate training objective and unweighted components, no optimizer."""
from __future__ import annotations
from dataclasses import dataclass
import torch
from .config import AblationConfig, finite_number
from .focal import occurrence_numerator
from .pinball import conditional_pinball_numerator
from .validation import finite_tensor, supervision


def combine_numerators(s_occ: torch.Tensor, s_qr: torch.Tensor, n_valid: int,
                       *, lambda_q: float) -> torch.Tensor:
    weight = finite_number("lambda_q", lambda_q)
    if weight <= 0 or type(n_valid) is not int or n_valid <= 0:
        raise ValueError("Positive lambda_q and positive integer N_valid required")
    for name, value in (("S_occ", s_occ), ("S_qr", s_qr)):
        if not isinstance(value, torch.Tensor) or value.ndim != 0 or value.layout != torch.strided or value.dtype not in (torch.float32, torch.float64):
            raise ValueError(f"{name}: FP32/FP64 scalar Tensor required")
    if s_occ.device != s_qr.device:
        raise ValueError("Numerator device mismatch")
    if s_occ.device.type not in ("cpu", "cuda"):
        raise ValueError("Numerator device must be CPU/CUDA")
    for name, value in (("S_occ", s_occ), ("S_qr", s_qr)):
        finite_tensor(name, value)
        if bool(value < 0):
            raise ValueError(f"{name}: loss numerator must be nonnegative")
    result = (s_occ + weight * s_qr) / n_valid  # weight exactly once
    finite_tensor("weighted training objective", result)
    return result


@dataclass(frozen=True)
class CandidateLossResult:
    experiment_id: str
    s_occ: torch.Tensor
    s_qr: torch.Tensor
    training_objective: torch.Tensor
    weighted_quantile_per_valid: torch.Tensor
    n_valid: int
    n_rain: int
    conditional_skipped: bool
    status: str = "VALID_SYNTHETIC_OR_CANDIDATE_LOSS_CALL"
    execution_authorization: bool = False

    @property
    def scientific_conditional_pinball(self) -> torch.Tensor | None:
        # No lambda here. None explicitly means not estimable, not zero skill.
        return self.s_qr / self.n_rain if self.n_rain else None


def candidate_loss(logit: torch.Tensor, qlog: torch.Tensor, rate: torch.Tensor,
                   imerg_valid: torch.Tensor, yunnan_mask: torch.Tensor,
                   *, config: AblationConfig) -> CandidateLossResult:
    if type(config) is not AblationConfig:
        raise ValueError("Explicit closed AblationConfig required")
    config.__post_init__()  # revalidate even if external code bypassed frozen setattr
    target = supervision(logit, qlog, rate, imerg_valid, yunnan_mask)
    s_occ = occurrence_numerator(logit, target.rainy, target.valid, alpha=config.alpha, gamma=config.gamma)
    with torch.autocast(logit.device.type, enabled=False):
        s_qr = conditional_pinball_numerator(qlog, torch.log1p(target.clean_rate.double()), target.rainy_valid)
        total = combine_numerators(s_occ, s_qr, target.n_valid, lambda_q=config.lambda_q)
        weighted = config.lambda_q * s_qr / target.n_valid
    return CandidateLossResult(config.experiment_id, s_occ, s_qr, total, weighted,
                               target.n_valid, target.n_rain, target.n_rain == 0)
\end{lstlisting}


\subsection{metrics.py：实际源码、函数位置与副本}

来源 src/\allowbreak{}yuntapr/\allowbreak{}experimental/\allowbreak{}phase\_\allowbreak{}b\_\allowbreak{}v2\_\allowbreak{}ablations/\allowbreak{}metrics.py。

SHA256：3450edb27553\allowbreak{}39801a5fd558\allowbreak{}64091b1c22e0\allowbreak{}628fa49fce71\allowbreak{}bc5bff72811f\allowbreak{}9021。参数语义见对应walkthrough；下面是实际完整注释源码副本。

\begingroup\small\setlength{\tabcolsep}{3pt}\renewcommand{\arraystretch}{1.12}
\begin{longtable}{>{\raggedright\arraybackslash}p{\dimexpr(\linewidth-6\tabcolsep)/3\relax}>{\raggedright\arraybackslash}p{\dimexpr(\linewidth-6\tabcolsep)/3\relax}>{\raggedright\arraybackslash}p{\dimexpr(\linewidth-6\tabcolsep)/3\relax}}
\toprule
函数或方法（参数名） & 源码行 & 阅读入口 \\
\midrule\endfirsthead
函数或方法（参数名） & 源码行 & 阅读入口 \\
\midrule\endhead
common\_\allowbreak{}validation\_\allowbreak{}sums(logit, qlog, rate, imerg\_\allowbreak{}valid, yunnan\_\allowbreak{}mask) & 38 & 完整type/\allowbreak{}default见下面源码 \\
\_\allowbreak{}\_\allowbreak{}post\_\allowbreak{}init\_\allowbreak{}\_\allowbreak{}(self) & 17 & 完整type/\allowbreak{}default见下面源码 \\
from\_\allowbreak{}candidate(cls, result) & 24 & 完整type/\allowbreak{}default见下面源码 \\
conditional\_\allowbreak{}pinball(self) & 28 & 完整type/\allowbreak{}default见下面源码 \\
pooled(self, other) & 31 & 完整type/\allowbreak{}default见下面源码 \\
\bottomrule\end{longtable}\endgroup

\begin{lstlisting}
"""Unweighted reporting interface; common core always uses E0 in FP64."""
from __future__ import annotations
from dataclasses import dataclass
import math
import torch
from .config import get_config
from .total import CandidateLossResult, candidate_loss


@dataclass(frozen=True)
class ReportingSums:
    s_occ: float
    s_qr: float
    n_valid: int
    n_rain: int

    def __post_init__(self) -> None:
        if not all(type(x) in (float, int) and math.isfinite(x) and x >= 0 for x in (self.s_occ, self.s_qr)):
            raise ValueError("Finite nonnegative unweighted reporting sums required")
        if type(self.n_valid) is not int or type(self.n_rain) is not int or not 0 <= self.n_rain <= self.n_valid or self.n_valid <= 0:
            raise ValueError("Positive N_valid and integer 0<=N_rain<=N_valid required")

    @classmethod
    def from_candidate(cls, result: CandidateLossResult) -> ReportingSums:
        return cls(float(result.s_occ.detach()), float(result.s_qr.detach()), result.n_valid, result.n_rain)

    @property
    def conditional_pinball(self) -> float | None:
        return self.s_qr / self.n_rain if self.n_rain else None

    def pooled(self, other: ReportingSums) -> ReportingSums:
        # Pool sums/counts rather than averaging unequal-batch metric values.
        return ReportingSums(self.s_occ + other.s_occ, self.s_qr + other.s_qr,
                             self.n_valid + other.n_valid, self.n_rain + other.n_rain)


@torch.no_grad()
def common_validation_sums(logit: torch.Tensor, qlog: torch.Tensor, rate: torch.Tensor,
                           imerg_valid: torch.Tensor, yunnan_mask: torch.Tensor) -> ReportingSums:
    """Fixed gamma2/lambda1 arithmetic, not an arm's weighted train objective.

    Caller-supplied tensors only. This function does not run a model or read data.
    The current tests call it exclusively with explicitly synthetic tensors.
    """
    result = candidate_loss(logit.double(), qlog, rate, imerg_valid, yunnan_mask, config=get_config("E0"))
    return ReportingSums.from_candidate(result)
\end{lstlisting}


\subsection{readiness.py：实际源码、函数位置与副本}

来源 src/\allowbreak{}yuntapr/\allowbreak{}experimental/\allowbreak{}phase\_\allowbreak{}b\_\allowbreak{}v2\_\allowbreak{}ablations/\allowbreak{}readiness.py。

SHA256：2f6ca9f1f5cf\allowbreak{}0da458cc008e\allowbreak{}05bb24d4c1ec\allowbreak{}d776172f35bf\allowbreak{}1b8cef81f234\allowbreak{}92e3。参数语义见对应walkthrough；下面是实际完整注释源码副本。

\begingroup\small\setlength{\tabcolsep}{3pt}\renewcommand{\arraystretch}{1.12}
\begin{longtable}{>{\raggedright\arraybackslash}p{\dimexpr(\linewidth-6\tabcolsep)/3\relax}>{\raggedright\arraybackslash}p{\dimexpr(\linewidth-6\tabcolsep)/3\relax}>{\raggedright\arraybackslash}p{\dimexpr(\linewidth-6\tabcolsep)/3\relax}}
\toprule
函数或方法（参数名） & 源码行 & 阅读入口 \\
\midrule\endfirsthead
函数或方法（参数名） & 源码行 & 阅读入口 \\
\midrule\endhead
require\_\allowbreak{}sha(value, name) & 17 & 完整type/\allowbreak{}default见下面源码 \\
state\_\allowbreak{}metadata\_\allowbreak{}digest(tensors) & 36 & 完整type/\allowbreak{}default见下面源码 \\
audit\_\allowbreak{}fresh\_\allowbreak{}metadata(records) & 68 & 完整type/\allowbreak{}default见下面源码 \\
review\_\allowbreak{}checkpoint\_\allowbreak{}binding(metadata, expected, *, resume\_\allowbreak{}approval\_\allowbreak{}reference) & 99 & 完整type/\allowbreak{}default见下面源码 \\
\_\allowbreak{}\_\allowbreak{}post\_\allowbreak{}init\_\allowbreak{}\_\allowbreak{}(self) & 28 & 完整type/\allowbreak{}default见下面源码 \\
validate(self) & 56 & 完整type/\allowbreak{}default见下面源码 \\
describe(self) & 140 & 完整type/\allowbreak{}default见下面源码 \\
run(self, *args, **kwargs) & 147 & 完整type/\allowbreak{}default见下面源码 \\
\bottomrule\end{longtable}\endgroup

\begin{lstlisting}
"""Metadata-only identity checks and an unconditionally blocked runner stub.

No torch import, tensor state loading, dataset access or approval issuance.
Synthetically valid metadata never implies that real initial states exist.
"""
from __future__ import annotations
from dataclasses import asdict, dataclass
import hashlib
import json
import re
from typing import NoReturn, Sequence
from .config import RunSpec, proposed_runs, workload

INPUT_DIFFERENCES = {"backbone.enc0.conv1.weight", "backbone.enc0.skip.weight"}


def require_sha(value: object, name: str) -> None:
    if type(value) is not str or re.fullmatch(r"[0-9a-f]{64}", value) is None:
        raise ValueError(f"{name}: lowercase SHA256 required, never a path or credential")


@dataclass(frozen=True)
class TensorIdentity:
    shape: tuple[int, ...]
    dtype: str
    sha256: str

    def __post_init__(self) -> None:
        if type(self.shape) is not tuple or not self.shape or not all(type(n) is int and n > 0 for n in self.shape):
            raise ValueError("Nonempty positive integer shape required")
        if self.dtype != "float32":
            raise ValueError("Frozen candidate parameter dtype is float32")
        require_sha(self.sha256, "tensor identity")


def state_metadata_digest(tensors: dict[str, TensorIdentity]) -> str:
    if type(tensors) is not dict or not tensors or not all(type(k) is str and k and
            type(v) is TensorIdentity for k, v in tensors.items()):
        raise ValueError("Named tensor identity metadata required")
    for value in tensors.values():
        value.__post_init__()
    body = json.dumps({k: asdict(v) for k, v in sorted(tensors.items())}, sort_keys=True, separators=(",", ":"))
    return hashlib.sha256(body.encode()).hexdigest()


@dataclass(frozen=True)
class FreshRecord:
    run: RunSpec
    tensors: dict[str, TensorIdentity]
    state_sha256: str
    post_pair_rng_sha256: str
    epoch_order_sha256: str
    initializer_code_sha256: str
    historical_checkpoint_loaded: bool = False

    def validate(self) -> None:
        if type(self.run) is not RunSpec:
            raise ValueError("Closed candidate RunSpec required")
        self.run.__post_init__()
        for name in ("state_sha256", "post_pair_rng_sha256", "epoch_order_sha256", "initializer_code_sha256"):
            require_sha(getattr(self, name), name)
        if self.historical_checkpoint_loaded is not False:
            raise ValueError("Historical model/optimizer/RNG state transfer forbidden")
        if state_metadata_digest(self.tensors) != self.state_sha256:
            raise ValueError("Declared state digest differs from supplied tensor metadata")


def audit_fresh_metadata(records: Sequence[FreshRecord]) -> dict:
    """Review declarations only; cannot prove origins or instantiate a model."""
    if len(records) != 18 or not all(type(r) is FreshRecord for r in records):
        raise ValueError("Exactly 18 fresh metadata records required")
    for r in records:
        r.validate()
    table = {r.run.run_id: r for r in records}
    if set(table) != {r.run_id for r in proposed_runs()}:
        raise ValueError("Missing, duplicate or unregistered run identity")
    if len({r.initializer_code_sha256 for r in records}) != 1:
        raise ValueError("Initializer code identity differs across arms")
    for seed in (2026, 2027, 2028):
        subset = [r for r in records if r.run.seed == seed]
        if len({r.post_pair_rng_sha256 for r in subset}) != 1 or len({r.epoch_order_sha256 for r in subset}) != 1:
            raise ValueError("Same-seed RNG/order metadata differs")
        for model in ("B0_MATCHED_V2", "B1_V2"):
            arms = [r for r in subset if r.run.model == model]
            if len({r.state_sha256 for r in arms}) != 1:
                raise ValueError("Same-model same-seed full fresh states differ across arms")
        a = table[RunSpec("E0", "B0_MATCHED_V2", seed).run_id].tensors
        b = table[RunSpec("E0", "B1_V2", seed).run_id].tensors
        if a.keys() != b.keys():
            raise ValueError("B0/B1 parameter names differ")
        different = {k for k in a if a[k].shape != b[k].shape}
        if different != INPUT_DIFFERENCES or any(a[k] != b[k] for k in a if k not in different):
            raise ValueError("Only two frozen input-shape differences permitted; shared tensors must match")
    return {"metadata_consistent": True, "actual_model_instantiated": False,
            "actual_initialization_proven": False, "scope": "DECLARED_OR_SYNTHETIC_METADATA_ONLY",
            "can_launch_formal_training": False}


def review_checkpoint_binding(metadata: dict, expected: dict, *, resume_approval_reference: str | None = None) -> dict:
    """No files are opened and no state can be applied by this review.

    A supplied reference is not authenticated here and cannot authorize resume.
    Future formal integration must bind a real independent event to LAST SHA.
    """
    fields = {"run_id", "last_sha256", "model_sha256", "optimizer_sha256", "rng_sha256",
              "scheduler_sha256", "code_sha256", "protocol_sha256", "data_sha256", "init_sha256",
              "completed_epoch", "retained_updates", "boundary", "original_execution_reference"}
    if type(metadata) is not dict or set(metadata) != fields or type(expected) is not dict:
        raise ValueError("Closed checkpoint metadata required")
    if metadata["run_id"] not in {r.run_id for r in proposed_runs()}:
        raise ValueError("Checkpoint run identity is not in candidate matrix")
    for key in fields:
        if key.endswith("sha256"):
            require_sha(metadata[key], key)
    epoch, updates = metadata["completed_epoch"], metadata["retained_updates"]
    if type(epoch) is not int or not 1 <= epoch <= 9 or type(updates) is not int or updates != epoch * 5228:
        raise ValueError("Complete-epoch retained budget mismatch")
    if metadata["boundary"] != "COMPLETED_TRAIN_AND_VALIDATION_EPOCH_ONLY":
        raise ValueError("Partial epoch is not a legal LAST")
    if type(metadata["original_execution_reference"]) is not str or not metadata["original_execution_reference"]:
        raise ValueError("Original execution provenance reference missing")
    binding = {"run_id", "last_sha256", "code_sha256", "protocol_sha256", "data_sha256", "init_sha256"}
    if set(expected) != binding or any(metadata[k] != expected[k] for k in binding):
        raise ValueError("Expected LAST/code/protocol/data/init/run binding differs")
    if resume_approval_reference is not None and (type(resume_approval_reference) is not str or not resume_approval_reference):
        raise ValueError("Nonempty independent-event reference or None required")
    return {"metadata_consistent": True, "remaining_retained_updates": 47052 - updates,
            "next_epoch": epoch + 1 if epoch < 9 else None, "approval_authenticity_verified": False,
            "resume_approval_reference_present": resume_approval_reference is not None,
            "can_apply_checkpoint_state": False, "can_launch_formal_training": False}


class FormalExecutionBlocked(RuntimeError):
    pass


class BlockedRunner:
    """Design skeleton with no execution implementation, even after flag edits."""

    def describe(self) -> dict:
        return {"runs": [r.run_id for r in proposed_runs()], "budget": workload(),
                "runner_implemented": False, "can_launch_formal_training": False,
                "required": ["Independent scientific approval", "Researcher code review",
                             "Explicit formal integration approval", "Data and compute permissions",
                             "Real fresh/data/code identities", "Separately bound stage execution authorization"]}

    def run(self, *args: object, **kwargs: object) -> NoReturn:
        raise FormalExecutionBlocked("Candidate skeleton has NO TRAINING CAPABILITY; new reviewed runner and independent approvals required")
\end{lstlisting}
\clearpage\section{关键数学、解析梯度与源码对应}
\(v=m_{\rm IMERG}\land m_{\rm Yunnan}\)，\(z=\mathbf1\{y_{32}>\operatorname{float32}(0.1)\}\)，
\(N_v=\sum v\)，\(N_r=\sum vz\)；比较在提升精度前完成。
\begin{align*}
b(l,z)&=\max(l,0)-lz+\log(1+e^{-|l|}),&p_t&=e^{-b},\\
\alpha_z&=\alpha z+(1-\alpha)(1-z),&S_o&=\sum v\alpha_z(1-p_t)^\gamma b,\\
\tau_i&=(i-0.5)/32,&e_i&=\log(1+y)-q_i,\\
\rho_{\tau_i}(e_i)&=\max(\tau_i e_i,(\tau_i-1)e_i),&S_q&=\sum vz\,\frac1{32}\sum_i\rho_{\tau_i}(e_i),\\
L_{\rm train}&=(S_o+\lambda_q S_q)/N_v,&\operatorname{CPB}&=S_q/N_r.
\end{align*}
Focal对应focal.py中的BCEWithLogits、exp、pow、sum；Pinball对应pinball.py中的固定tau、maximum、mean、sum；total.py只乘一次lambda。\(\gamma=0,\alpha=.5\)仍为\(.5\,\mathrm{BCE}\)。CPB不带权重，无雨不可估计。
\[
\frac{\partial S_o}{\partial l}=\alpha_z(p-z)\{m^\gamma+\gamma b p_t m^{\gamma-1}\},\quad m=1-p_t,\quad\gamma>0;
\qquad\gamma=0:\ \alpha_z(p-z).
\]
\[
\frac{\partial L}{\partial q_i}=\frac{\lambda_q}{32N_v}\begin{cases}-\tau_i,&e_i>0,\\1-\tau_i,&e_i<0.\end{cases}
\]
以上梯度仅在有效位置并按相应雨标签生效；\(e=0\)为折点。test\_gradients.py独立核对解析式、FP64差分和全部32tau；不是实际模型性能结论。
\[
\eta(u)=\begin{cases}10^{-4}u/5228,&1\le u\le5228,\\
10^{-6}+(10^{-4}-10^{-6})\dfrac{1+\cos(\pi(u-5228)/(261400-5228))}{2},&5228<u\le47052.\end{cases}
\]
config.py的learning\_rate\_prefix逐位置与冻结纯数学函数对照；没有optimizer或正式更新。BlockedRunner只抛错。V2\_PHASE\_B\_AUTHORIZED=false，历史恢复NOT\_GRANTED，2025访问0。工程交付到此停止。
\end{document}
